% ============================================================================
% Structural limits of parcel-level deep reinforcement learning for cadastral
% layout optimization (working title in v5 — softened in v6 per reviewer)
% ============================================================================
% Target journal: Computers, Environment and Urban Systems
% CEUS initial submission package: anonymous main manuscript copy
% Editable LaTeX source retained for review.
% Version: v3 (double-anonymous review version)
% ============================================================================
\documentclass[]{interact}

\usepackage{epstopdf}
\usepackage{subfigure}
\usepackage{natbib}
\usepackage{amsmath,amssymb}
\usepackage{booktabs}
\usepackage{multirow}
\usepackage{graphicx}
\usepackage{algorithm}
\usepackage{algorithmic}
\usepackage{hyperref}

\theoremstyle{plain}
\newtheorem{theorem}{Theorem}[section]

\begin{document}

\articletype{RESEARCH ARTICLE}

\title{Diagnostic benchmarking of deep reinforcement learning for cadastral land-use decision support}

\author{
\name{Anonymous Authors}
}

\maketitle

\begin{abstract}
Computer-based land-use decision support increasingly considers deep reinforcement learning (DRL) for spatial layout optimisation, but practical deployment requires evidence that learned sequential policies add value over transparent local baselines. We test this question for cadastral-style farmland--forest exchange using a public synthetic-parcel benchmark generated from Copernicus DEM, Impact Observatory land-cover, and SLIC tessellation, heuristically matched to aggregate target ranges for mountainous cadastral settings. The evaluation contains 30 Maskable PPO policies across three regions (3{,}607--11{,}786 parcels), two reward-weighting strategies, and five seeds; Region B policies were retrained under the area-weighted environment used for evaluation. A slope-exact local baseline reduces area-weighted farmland slope by 5.9--9.7 percentage points across the three regions, but degrades contiguity. A seven-weight contiguity-aware exact-delta sweep recovers positive-contiguity local points in every region, dominates all six DRL mean outcomes, and dominates 27 of 30 individual deterministic DRL rollouts; the three undominated DRL rollouts prevent a universal dominance claim. A no-reuse action-mask parity audit leaves the slope-exact local outcomes unchanged and still dominates all six DRL mean outcomes. A top-5 two-step local lookahead changes one-step local slope outcomes by at most 0.154 percentage points, indicating that the local result is not an artefact of purely myopic ranking. Fixed-policy decoding at 25--200 paired exchanges and a smaller-budget prefix-horizon audit do not remove the slope-exact local advantage. Stochastic-vs-deterministic inference and neighbour-overlap analyses show that DRL learned spatial clustering, but weak inter-swap coupling and train--evaluation mismatch limit its deterministic decision-support value. These results do not support a universal claim about DRL. They show that, in the tested low-intervention and sparse-adjacency synthetic parcel regimes, Maskable PPO with a scorer-MLP policy did not justify its training cost relative to exact-delta local diagnostics. The benchmark supports a practical rule for CEUS-style decision support: run slope-only, slope-exact, and multi-objective local baselines before deploying GPU-intensive DRL.
\end{abstract}

\begin{keywords}
deep reinforcement learning; spatial optimisation; cadastral parcels; algorithm--structure mismatch; synthetic data benchmarking; near-independence
\end{keywords}

% ============================================================================
\section{Introduction}\label{sec:intro}
% ============================================================================

Computer-based land-use decision support must increasingly choose between transparent heuristics and data-intensive learning systems. In farmland layout planning, one recurring task is to rearrange which parcels are cropped versus afforested so that aggregate farmland slope is reduced while spatial contiguity is not unduly damaged. This task is relevant in mountainous regions where fragmented holdings create inefficiencies for agricultural production and ecological services \citep{liu2014land,long2018land}. With over 100 million hectares under such pressure in China alone \citep{wu2005consolidation,tan2006fragmentation}, computational methods that are reproducible, interpretable, and reliable under parcel-level constraints have practical value. The decision-support problem is not only to produce one allocation, but to decide whether a costly learned policy is justified when transparent local alternatives can be evaluated quickly.

Two algorithmic traditions dominate this problem space. \emph{Mathematical-programming and metaheuristic} approaches, including linear programming \citep{aerts2003using}, multi-objective evolutionary algorithms \citep{deb2002fast,stewart2004multiobjective,cao2011sustainable}, and simulated annealing \citep{cao2012spatial}, treat the optimisation as a static assignment or search problem. \emph{Deep reinforcement learning} has emerged as a sequential alternative for combinatorial optimisation \citep{bello2017neural,nazari2018reinforcement,kool2019attention,zheng2023spatial}, motivated by its ability to learn dependencies between successive decisions. For land-use decision support, the relevant question is not whether DRL can solve abstract allocation problems, but whether its additional training cost and operational complexity are justified on irregular parcel systems where simple local rankings may already capture much of the objective.

\subsection{The synthetic--real gap}

Published DRL studies in routing, allocation, and grid-like spatial planning often operate on regular abstractions with uniform unit size, dense neighbourhoods, or high intervention ratios \citep{bello2017neural,nazari2018reinforcement,kool2019attention,zheng2023spatial}. Cadastral-style parcels create a different computational regime: parcel areas vary, adjacency is sparse and irregular, and planning budgets usually affect only a small fraction of available units. Whether DRL's sequential-decision advantage survives this transition is therefore an empirical question. Two hypotheses bracket potential failure modes:

\begin{itemize}
\item \textbf{Scale hypothesis}: failure correlates with the number of swappable parcels because the action space grows combinatorially, exhausting the exploration budget;
\item \textbf{Structure hypothesis}: failure correlates with cadastral irregularity because heterogeneous parcel areas and sparse spatial interactions decouple sequential decisions, making greedy selection near-optimal regardless of scale.
\end{itemize}

Distinguishing these hypotheses is methodologically difficult because real cadastral data confounds scale, topography, parcel-area distribution, adjacency structure, and farm--forest ratio. Restricted-data accessibility further hampers reproducibility: Chinese parcel-level cadastral boundary data are not publicly accessible for open benchmarking. This creates a decision-support problem as much as a modelling problem. Without reproducible stress tests, practitioners cannot know whether a learned policy is adding planning value or merely approximating a transparent local heuristic at much higher cost.

\subsection{Approach}

This paper addresses the gap by building a public synthetic-parcel pipeline heuristically matched to aggregate structural ranges reported for mountainous cadastral settings. The pipeline uses Copernicus DEM (30 m), Impact Observatory land-cover \citep{karra2021global} (10 m, year 2019), and SLIC superpixel tessellation \citep{achanta2012slic}; no restricted cadastral boundaries are consumed. Three synthetic regions (3{,}607--11{,}786 parcels) are subjected to a fully specified Maskable Proximal Policy Optimisation pipeline: an MDP formulation with enforced paired action masking, a permutation-invariant scoring policy network, two reward-weighting strategies, and area-weighted slope/contiguity objectives. The final reported DRL evaluation contains 30 trained policies (three regions $\times$ two configurations $\times$ five seeds), with the Region B policies retrained under the area-weighted environment used for evaluation. These policies are compared against Random, slope-only Greedy, slope-exact Greedy, a seven-weight contiguity-aware exact-delta local sweep, a no-reuse action-mask parity audit of that sweep, a top-5 two-step lookahead variant of the local sweep, and Genetic Algorithm baselines. A stochastic-versus-deterministic inference comparison and a neighbour-overlap analysis then diagnose why the learned policy fails to match local ranking.

\subsection{Contributions}

\begin{enumerate}
\item A reproducible synthetic-parcel pipeline heuristically matched to published aggregate structural ranges from public sources alone (Section~\ref{sec:synth}). The pipeline parametrises CV(area), average degree, and farm:forest ratio independently, enabling controlled stress tests not possible on a single restricted cadastral sample.

\item A fully specified parcel-level DRL benchmark for cadastral layout optimisation (Section~\ref{sec:methods}), including MDP formulation, enforced paired action masking, a permutation-invariant scoring policy network, Lagrangian dual-variable adaptation of reward weights, and an area-weighted training/evaluation environment.

\item 30 paired-inference evaluations spanning three regions and two configurations (Section~\ref{sec:results}). Under area-weighted evaluation, a slope-exact Greedy baseline outperforms DRL on slope across all three regions but substantially degrades contiguity.

\item A contiguity-aware exact-delta local sweep over seven scalarisation weights (Section~\ref{sec:results-pareto}). The sampled local Pareto curve contains points that dominate all six DRL mean outcomes and 27 of 30 individual deterministic DRL runs, while leaving three strong individual DRL runs undominated.

\item A no-reuse action-mask parity audit (Section~\ref{sec:results-mask-parity}) testing whether the local result depends on allowing parcels to re-enter the feasible set after conversion. The slope-exact local outcomes are unchanged, and the no-reuse local grid still dominates all six DRL mean outcomes.

\item A top-5 two-step local lookahead audit (Section~\ref{sec:results-lookahead}) testing whether one-step exact-delta ranking misses sequential gains. The audit changes local slope outcomes by at most 0.154 percentage points across the tested region--weight combinations.

\item A decision-support diagnostic workflow (Section~\ref{sec:disc-implications}) that tells practitioners when to run inexpensive local baselines before committing to GPU-intensive DRL training.

\item Direct empirical evidence that train-evaluation distribution mismatch is a partial driver of the failure (Section~\ref{sec:results-stochastic}): stochastic-sampling paired-inference improves Region B's slope outcome by approximately 1.3 percentage points relative to deterministic argmax, but does not close the gap to slope-only, slope-exact, or contiguity-aware local ranking.

\item An empirical neighbour-overlap analysis (Section~\ref{sec:results-overlap}) quantifying whether learned policies create spatially coupled exchange sequences under the tested synthetic-parcel settings.

\item An empirical distinction between bimodal seed convergence---a structural property of the optimisation landscape---and reward-weighting failures, demonstrated by the same RNG seed failing under both Lagrangian and fixed-weight configurations (Section~\ref{sec:results-bimodal}).
\end{enumerate}

% ============================================================================
\section{Synthetic cadastral data generation}\label{sec:synth}
% ============================================================================

\subsection{Design rationale}

Real Chinese cadastral data from the TNLS database is access-restricted, foreclosing public reproducibility. To enable open benchmarking while approximating structural properties that define the cadastral domain, this study constructs a synthetic pipeline targeting aggregate ranges consistent with TNLS-like mountainous cadastral settings reported in the literature \citep{wu2005consolidation,tan2006fragmentation,niu2023ffi}. The pipeline reads only public products and releases all intermediates under CC-BY-4.0.

\subsection{Public data sources}

Two raster products provide the topographic and land-cover signal:
\begin{itemize}
\item \textbf{Copernicus DEM GLO-30} \citep{esa2021copernicusdem} (30~m resolution): provides per-pixel elevation; slope and aspect are derived using Horn's method \citep{horn1981slope}. Accessed via Google Earth Engine asset \texttt{COPERNICUS/DEM/GLO30}.
\item \textbf{Impact Observatory annual land-cover} \citep{karra2021global}, year 2019 to align with the TNLS reference year, 10~m resolution, 9-class scheme. Class codes are: 1 (water), 2 (trees), 4 (flooded vegetation), 5 (crops), 7 (built), 8 (bare), 9 (snow), 10 (clouds), 11 (rangeland). Accessed via the GEE community catalogue asset \texttt{projects/sat-io/open-datasets/landcover/ESRI\_Global-LULC\_10m\_TS}.
\end{itemize}

A geographic extent of $\sim$50$\times$50~km in transitional Sichuan Basin terrain (lon $105.7$--$106.4$, lat $29.2$--$29.9$) is selected to provide topographic heterogeneity comparable to the target mountainous setting.

\subsection{Tessellation and feature assignment}

Cadastral-like polygons are produced by SLIC superpixel segmentation \citep{achanta2012slic} on a stacked feature image combining DEM, slope, and land-cover. The number of segments is over-requested at $2N$ where $N$ is the target parcel count; the resulting partition is trimmed by stratified random drop preferring the smallest fragments. Three target sizes are generated: $N \in \{3{,}607,\, 5{,}500,\, 11{,}786\}$. The third target is an approximate match for $\sim$12{,}950 parcels; the achieved 11{,}786 ($\sim$91\% of target) represents the largest synthetic region for which SLIC tessellation produced a topologically clean partition under the chosen extent.

For each parcel, ten features are computed (matching a standard parcel-feature schema): land-use type (one-hot, 3 dims), normalised slope, normalised area, proportion of same-type neighbours, neighbour count, mean neighbour slope, normalised elevation, aspect sine, aspect cosine, and shape index ($P / 2\sqrt{\pi A}$). Slope and aspect are derived from the DEM via the Horn finite-difference operator \citep{horn1981slope}. Land-use type is assigned by ranking parcels by their LULC overlay scores, taking the top-$N_{\text{farm}}$ farmland-scoring parcels and the top-$N_{\text{forest}}$ forest-scoring among the remainder. Adjacency is constructed via Queen contiguity using \texttt{libpysal} \citep{rey2007pysal}.

\subsection{Calibration verification}

Table~\ref{tab:calibration} compares synthetic regions against published TNLS-like aggregate statistics that serve as \emph{heuristic design targets} for the pipeline. We do not claim that each target value is independently calibrated from a single exact published source---only that the values are plausible given the literature on mountainous Chinese cadastral structure \citep{wu2005consolidation,tan2006fragmentation,niu2023ffi,long2018land}. Four of five structural targets are achieved across all three regions; CV(area) is below the heuristic target, a known limitation of SLIC tessellation discussed in Section~\ref{sec:disc-limits}.

\begin{table}[htbp]
\centering
\caption{Synthetic cadastral data against heuristic design targets consistent with TNLS-like mountainous cadastral settings. CV(area) below target reflects SLIC tessellation's tendency toward area uniformity; this is acknowledged as a limitation (Section~\ref{sec:disc-limits}). Heuristic targets motivated by \citet{wu2005consolidation,tan2006fragmentation,niu2023ffi}; the target values are not claimed to come from a single calibration dataset but rather summarise typical ranges in the cited literature.}
\label{tab:calibration}
\begin{tabular}{lrrrr}
\toprule
Property & Target & Region A & Region B & Region C \\
\midrule
Parcel count           & 3{,}607 / 5{,}500 / 12{,}950$^\dagger$ & 3{,}607 \checkmark & 5{,}500 \checkmark & 11{,}786 ($\sim$) \\
Farmland count         & 1{,}130 / 1{,}998 / 5{,}032 & 1{,}130 \checkmark & 1{,}998 \checkmark & 5{,}032 \checkmark \\
Forest count           & 471 / 671 / 1{,}755 & 471 \checkmark & 671 \checkmark & 1{,}755 \checkmark \\
Farmland:forest ratio  & 2.40 / 2.98 / 2.87 & 2.40 \checkmark & 2.98 \checkmark & 2.87 \checkmark \\
Mean Queen degree      & $\sim$ 5.0 & 5.28 \checkmark & 4.47 \checkmark & 5.61 \checkmark \\
Std Queen degree       & $\sim$ 2.3 & 1.84 \checkmark & 1.57 \checkmark & 1.89 \checkmark \\
Mean farmland slope    & 8--10$^\circ$ & 7.68$^\circ$ \checkmark & 7.66$^\circ$ \checkmark & 6.61$^\circ$ \checkmark \\
Elevation range (m)    & 180--885 & 284--815 \checkmark & 246--974 \checkmark & 262--726 \checkmark \\
CV(parcel area)        & $\sim$ 2.0 & 0.85 (low) & 0.57 (low) & 0.79 (low) \\
\bottomrule
\multicolumn{5}{l}{\footnotesize $^\dagger$ Region C target is approximate; SLIC tessellation under the chosen extent}\\
\multicolumn{5}{l}{\footnotesize achieved 11{,}786 parcels ($\sim$91\% of target).}\\
\end{tabular}
\end{table}

\subsection{Pipeline algorithm}

Algorithm~\ref{alg:synth} summarises the synthetic-data construction. The pipeline is fully deterministic given a seed; complete code is released.

\begin{algorithm}[htbp]
\caption{Synthetic cadastral data construction.}
\label{alg:synth}
\begin{algorithmic}[1]
\REQUIRE Public extent $E$ in lon/lat; target counts $(N, N_\text{farm}, N_\text{forest})$; RNG seed.
\STATE Fetch DEM raster $R_\text{DEM}$ from Copernicus GLO-30 over $E$.
\STATE Fetch land-cover raster $R_\text{LC}$ from Impact Observatory 2019 over $E$.
\STATE Compute slope $R_S$ and aspect $R_A$ from $R_\text{DEM}$ via Horn's method.
\STATE Stack $(R_\text{DEM}, R_S, R_\text{LC})$ as a 3-channel feature image.
\STATE Run SLIC superpixel segmentation requesting $2N$ segments with compactness 8.
\STATE Polygonize the label raster; reproject to EPSG:4523 (CGCS2000 GK-35).
\STATE Stratified-random drop smallest fragments to retain exactly $N$ parcels.
\STATE For each parcel $p$, compute LULC overlay scores $\sigma_\text{farm}(p), \sigma_\text{forest}(p)$.
\STATE Assign $\text{type}(p) \leftarrow \text{farmland}$ for the top-$N_\text{farm}$ parcels by $\sigma_\text{farm}$ that prefer farmland over forest.
\STATE Assign $\text{type}(p) \leftarrow \text{forest}$ for the top-$N_\text{forest}$ remaining by $\sigma_\text{forest}$.
\STATE Build Queen contiguity graph using \texttt{libpysal}.
\STATE Compute per-parcel features: slope, aspect sin/cos, area, shape index, neighbour stats.
\STATE Output: (i) GeoDataFrame of parcels with attributes; (ii) sparse adjacency CSR.
\end{algorithmic}
\end{algorithm}

% ============================================================================
\section{Methods}\label{sec:methods}
% ============================================================================

\subsection{MDP formulation}

The land-use exchange problem is cast as a Markov Decision Process $(\mathcal{S}, \mathcal{A}, T, R, \gamma)$, fully specified below.

\textbf{State} $s_t \in \mathcal{S}$ is a structured tuple $(F_t, G_t, M_t)$, where $F_t \in \mathbb{R}^{N \times K}$ is a per-parcel feature matrix ($K=10$ features per parcel: land-use type one-hot, normalised slope, normalised area, proportion of same-type neighbours, neighbour count, mean neighbour slope, normalised elevation, aspect sine, aspect cosine, shape index), $G_t \in \mathbb{R}^8$ is an 8-dimensional global statistics vector (current farmland count, mean slope, mean contiguity, plus second-moment statistics), and $M_t \in \{0,1\}^N$ is a Boolean action mask indicating which parcels are valid choices at step $t$ (Section~\ref{sec:adapt}).

\textbf{Action} $a_t \in \{1, \ldots, N\}$ selects a single parcel from those allowed by $M_t$. Selection alternates between farmland-only (odd $t$) and forest-only (even $t$) under enforced paired action masking (Section~\ref{sec:adapt}).

\textbf{Transition} $T$ is deterministic: when $a_t$ is a farmland parcel, it is held in a buffer; on the next step's forest selection, the two parcels' land-use labels are swapped, producing one paired exchange. The episode budget is fixed at $B = 100$ paired exchanges (200 raw steps).

\textbf{Reward} $r_t$ is shaped to encode all three objectives:
\begin{equation}
r_t = w_{\text{slope}}\, \Delta \bar{s}_t + w_{\text{cont}}\, \Delta \bar{c}_t - w_{\text{count}} (\Delta n_t)^2 + b_{\text{pair}},
\label{eq:reward}
\end{equation}
where $\Delta \bar{s}_t = \bar{s}_t - \bar{s}_{t-1}$ is the change in area-weighted mean farmland slope, defined as $\bar{s}_t = \sum_{i \in F_t} s_i a_i / \sum_{i \in F_t} a_i$ for the farmland set $F_t$, parcel slope $s_i$, and area $a_i$ (negative is desirable). $\Delta \bar{c}_t = \bar{c}_t - \bar{c}_{t-1}$ is the change in mean farmland contiguity, defined as
\begin{equation}
\bar{c}_t = \frac{1}{|F_t|} \sum_{i \in F_t} |\{j \in N(i) : \mathrm{type}(j) = \mathrm{farmland}\}|,
\label{eq:contiguity}
\end{equation}
the average number of farmland Queen neighbours per farmland parcel (positive is desirable). $\Delta n_t$ is the deviation from initial farmland count, and $b_{\text{pair}}$ is a fixed bonus awarded only after a complete farmland--forest swap pair. The reward weights $w_{\text{slope}}, w_{\text{cont}}, w_{\text{count}}$ are configuration-specific: under the EntOnly configuration they are static at $w_{\text{slope}}=1000$, $w_{\text{cont}}=500$, $w_{\text{count}}=1000$; under the Lagrangian configuration $w_{\text{cont}}$ and $w_{\text{count}}$ are dynamic dual variables (Section~\ref{sec:adapt}). The discount factor $\gamma = 0.995$ encourages long-horizon planning.

\textbf{Title-of-objective vs evaluation metric.} The training reward in Eq.~\eqref{eq:reward} and the paired-inference evaluation metric (also area-weighted mean slope change) use the same area-weighted formulation throughout, ensuring that improvements observed during training-time stochastic rollouts are commensurate with paired-inference evaluation outcomes.

\textbf{Policy architecture.} The policy is a \emph{PlanningUnitScoringPolicy}: a parcel-permutation-invariant network that scores each parcel independently using a shared multi-layer perceptron over the per-parcel feature vector concatenated with the global statistics ($K + |G| = 18$ inputs), producing a scalar logit per parcel. The action distribution is the masked softmax of these logits restricted to mask-allowed parcels. Concretely, the scoring MLP has hidden layers $[128, 64]$ with ReLU activations and outputs a scalar; the value network is a separate MLP $[128, 64]$ over the 8 global statistics. Because the scoring network operates per parcel with shared weights, the same trained model can in principle be applied to regions of different parcel counts without architectural changes.

\subsection{Two methodological adaptations}\label{sec:adapt}

\paragraph{Enforced paired action masking.} Without it, the agent learns degenerate ``do-nothing'' policies because the count penalty term in \eqref{eq:reward} can be minimised by avoiding all swaps. Forcing the action mask to alternate between farmland (odd $t$) and forest (even $t$) selection pools eliminates this degenerate equilibrium and guarantees farmland-count conservation at the end of every paired step. The count penalty $w_{\text{count}}(\Delta n_t)^2$ remains non-zero at intermediate odd-$t$ steps where one farmland has been removed but the forest has not yet been added, so $\Delta n_t = -1$ before pair completion. Retaining the penalty---and adapting $w_{\text{count}}$ via the Lagrangian formulation---provides a within-pair signal that discourages selecting farmland parcels for which no suitable forest counterpart is available, complementing the mask which only enforces alternation but not coupling between the two selections within a pair.

\paragraph{Lagrangian dual-variable adaptation.} On real cadastral data the relative magnitudes of slope and contiguity rewards are unknown \emph{a priori}, motivating the Lagrangian formulation \citep{stooke2020responsive,tessler2019reward}:
\begin{align}
\mu_{\text{cont}}^{(k+1)} &= \mathrm{clip}\big(\mu_{\text{cont}}^{(k)} + \alpha_{\text{cont}} (d_{\text{cont}} - \bar{c}^{(k)}),\, 100,\, 2{,}000\big),\\
\mu_{\text{count}}^{(k+1)} &= \mathrm{clip}\big(\mu_{\text{count}}^{(k)} + \alpha_{\text{count}} (\bar{v}^{(k)} - d_{\text{count}}),\, 100,\, 2{,}000\big).
\end{align}
$w_{\text{slope}} = 1{,}000$ remains fixed as the primary objective. Entropy coefficient annealing from 0.05 to 0.005 over the first 50\% of training is applied to all configurations.

\subsection{Two configurations}

Two configurations are compared: \textbf{Lagrangian} (adaptive $\mu_{\text{cont}}, \mu_{\text{count}}$ with entropy annealing) and \textbf{EntOnly} (constant weights $w_{\text{slope}}=1000$, $w_{\text{cont}}=500$, $w_{\text{count}}=1000$ with entropy annealing only). Both configurations enforce paired action masking. Testing both dynamic (Lagrangian) and static (EntOnly) reward-weighting strategies is a robustness check: if both fail in the same regions, the failure cannot be attributed to a poor static hyper-parameter choice, and conversely if Lagrangian rescues a region where EntOnly fails, the failure is reward-weighting--specific. The results in Section~\ref{sec:results} show that both configurations produce qualitatively identical paired-inference outcomes, supporting the former interpretation.

\subsection{Baselines}

\textbf{Random}: pairs farmland and forest parcels uniformly at random, 5 seeds.\\
\textbf{Slope-only Greedy}: deterministically pairs the highest-slope farmland parcel with the lowest-slope forest parcel each step. Single deterministic run.\\
\textbf{Exact-delta Greedy}: at each step, evaluates the exact change in area-weighted mean farmland slope for every feasible (farmland $i$, forest $j$) pair and selects the pair minimising the new mean:
$$\bar{s}' = \frac{S_F - s_i a_i + s_j a_j}{A_F - a_i + a_j}.$$
At each of the $B = 100$ paired exchanges this requires evaluating $|F| \times |O|$ pairs (e.g.\ $5{,}032 \times 1{,}755 \approx 8.8 \cdot 10^6$ pairs in Region~C). The exhaustive search makes exact-delta Greedy a strong one-step local-ranking baseline for the slope objective. Single deterministic run.\\
\textbf{Contiguity-aware exact-delta local ranking}: for the same feasible pair set, evaluates both the exact area-weighted slope change and the exact contiguity change, then selects
\begin{equation}
(i^\ast,j^\ast) = \arg\max_{i \in F,\, j \in O}
\left[
\frac{\bar{s}_t-\bar{s}'_{ij}}{|\bar{s}_0|+\epsilon}
+ \lambda_{\mathrm{cont}}
\frac{\bar{c}'_{ij}-\bar{c}_t}{|\bar{c}_0|+\epsilon}
\right].
\label{eq:local-scalar}
\end{equation}
Here $F$ is the current farmland set, $O$ is the current forest set, $\epsilon=10^{-12}$ prevents division by zero, and $\lambda_{\mathrm{cont}}$ controls the slope--contiguity tradeoff. The contiguity delta is computed analytically. Let $q_k$ denote the current number of farmland neighbours of parcel $k$ and $Q=\sum_{k\in F}q_k$; then for removing farmland $i$ and adding forest $j$,
$$Q' = Q - 2q_i + 2(q_j-\mathbb{1}[j\in N(i)]), \qquad \bar{c}'_{ij}=Q'/|F|.$$
We sweep $\lambda_{\mathrm{cont}}\in\{0,0.25,0.5,1,2,4,8\}$. The $\lambda_{\mathrm{cont}}=0$ run exactly reproduces Exact-delta Greedy, providing an internal check on the implementation. Single deterministic run per weight and region.\\
\textbf{No-reuse local parity audit}: the DRL environment marks selected swappable parcels as converted, preventing the same parcel from being selected again within an episode. The main local sweep ranks over the current farmland and forest sets. To test whether this implementation difference affects the comparison, we re-run the seven-weight local sweep with an explicit no-reuse mask that removes both parcels in each committed swap from subsequent candidate sets. This audit matches the DRL episode-level converted-parcel constraint but remains a deterministic exact-delta diagnostic rather than a globally optimised planner.\\
\textbf{Two-step lookahead local ranking}: to test whether one-step exact-delta ranking is too myopic, we add a bounded lookahead diagnostic for $\lambda_{\mathrm{cont}}\in\{0,0.5,1\}$. At each paired exchange, the algorithm first selects the top five immediate scalarised swaps, temporarily applies each candidate, evaluates the best next immediate scalarised swap, and commits the first swap with the largest two-step score. This audit is a transparent two-step diagnostic rather than a globally optimal planner.\\
\textbf{Genetic Algorithm}: tournament selection (size 3), single-point crossover (rate 0.8), random mutation (0.1), population 50, 100 generations, 5 seeds. Fitness mirrors the DRL reward. GA is included as a conventional stochastic baseline drawn from the land-use optimisation literature \citep{cao2011sustainable,stewart2004multiobjective}, not as an exhaustively tuned optimiser.

\subsection{Training and paired-inference evaluation}

All DRL models train under MaskablePPO \citep{huang2022closer}---a masked variant of Proximal Policy Optimisation \citep{schulman2017proximal}---from the \texttt{sb3-contrib} extension of the Stable-Baselines3 library \citep{raffin2021stablebaselines}: $10^6$ timesteps, 2{,}048-step rollouts, mini-batch 256, learning rate $3{\times}10^{-4}$, $\gamma=0.995$, $\lambda_{\text{GAE}}=0.95$, clip range 0.2. The PlanningUnitScoringPolicy uses hidden layers $[128, 64]$. Training was conducted on NVIDIA A100 GPUs; per-region training time ranged from 1.5 to 3.5~hours.

Two configurations $\times$ three regions $\times$ five seeds = 30 trained models in total. The training-time rollout is stochastic: actions sampled from the masked softmax distribution. Evaluation uses a single deterministic episode of 100 paired exchanges where actions are taken as the $\arg\max$ over the masked logit distribution. Both phases use enforced paired action masking and farmland-count conservation; the asymmetry between training-time stochastic exploration and evaluation-time deterministic argmax is the source of the train--evaluation distribution mismatch reported in Section~\ref{sec:results-mismatch}---the mismatch is in action-distribution sharpness, not in masking. Reported metrics are slope and contiguity changes from this paired-inference run, not training-rollout averages.

% ============================================================================
\section{Results}\label{sec:results}
% ============================================================================

\subsection{Cross-region method comparison}\label{sec:results-main}

Table~\ref{tab:main} presents the primary paired-inference comparison across the three synthetic regions, evaluated under the area-weighted environment that matches Eq.~\eqref{eq:reward}. The convergence rate, defined as the fraction of seeds achieving slope reduction beyond $-1\%$, is reported for the stochastic methods (Random, GA, both DRL configurations). For Region B, the DRL models reported here were retrained from scratch under the area-weighted environment; Region A and C reuse the original 30-model sweep but are re-evaluated under the area-weighted environment.

Two slope-focused Greedy variants are reported in Table~\ref{tab:main}. \textbf{Slope-only Greedy} ranks farmland and forest parcels by their slopes alone and pairs them in extreme order---a standard heuristic baseline. \textbf{Exact-delta Greedy} computes the exact one-step change in the area-weighted mean-slope objective for every feasible swap pair and chooses the optimal one each step. Slope-only Greedy ignores parcel area, so under an area-weighted objective it is not generally optimal; the exact-delta variant is the strongest natural local-ranking baseline for the single slope objective. Table~\ref{tab:local-pareto} then tests whether adding contiguity explicitly to the same exact-delta local framework changes the DRL comparison.

\begin{table}[htbp]
\centering
\caption{Primary paired-inference evaluation across methods and regions under the area-weighted environment. Slope change: more negative is better; contiguity change: more positive is better. DRL means are reported with descriptive bootstrap 95\% intervals from five seeds (5{,}000 bootstrap samples); raw per-seed outcomes are reported in Appendix~\ref{app:per-seed}. The scalarised local sweep is reported separately in Table~\ref{tab:local-pareto}.}
\label{tab:main}
\small
\begin{tabular}{llrrrl}
\toprule
Region & Method & $n$ & Slope $\Delta$\,(\%) & Cont. $\Delta$ & Conv.\\
\midrule
\multirow{6}{*}{A (3{,}607 parcels)}
& Random                   & 5 & $+9.32 \pm 0.03$ & $-0.512 \pm 0.030$ & 0/5 \\
& GA                       & 5 & $+1.65 \pm 0.02$ & $-0.022 \pm 0.016$ & 0/5 \\
& Greedy (slope-only)      & 1 & $-0.56$ & $+0.147$ & --- \\
& \textbf{Greedy (exact-delta)} & 1 & $\mathbf{-9.66}$ & $-0.549$ & --- \\
& DRL-Lagrangian           & 5 & $+1.02$ \tiny{[$-0.21, +3.15$]} & $+0.112 \pm 0.065$ & 0/5 \\
& DRL-EntOnly              & 5 & $+0.63$ \tiny{[$-0.73, +2.93$]} & $+0.067 \pm 0.061$ & 0/5 \\
\midrule
\multirow{6}{*}{B (5{,}500 parcels)}
& Random                   & 5 & $+5.60 \pm 0.04$ & $-0.215 \pm 0.011$ & 0/5 \\
& GA                       & 5 & $-1.44 \pm 0.01$ & $-0.010 \pm 0.005$ & 5/5 \\
& Greedy (slope-only)      & 1 & $-4.30$ & $+0.035$ & --- \\
& \textbf{Greedy (exact-delta)} & 1 & $\mathbf{-7.64}$ & $-0.174$ & --- \\
& DRL-Lagrangian           & 5 & $+2.19$ \tiny{[$+2.03, +2.31$]} & $-0.164 \pm 0.003$ & 0/5 \\
& DRL-EntOnly              & 5 & $+2.16$ \tiny{[$+1.86, +2.35$]} & $-0.127 \pm 0.081$ & 0/5 \\
\midrule
\multirow{6}{*}{C (11{,}786 parcels)}
& Random                   & 5 & $+3.79 \pm 0.01$ & $-0.138 \pm 0.008$ & 0/5 \\
& GA                       & 5 & $-0.24 \pm 0.01$ & $+0.022 \pm 0.007$ & 0/5 \\
& Greedy (slope-only)      & 1 & $-1.88$ & $+0.091$ & --- \\
& \textbf{Greedy (exact-delta)} & 1 & $\mathbf{-5.87}$ & $-0.161$ & --- \\
& DRL-Lagrangian           & 5 & $-2.20$ \tiny{[$-2.75, -1.45$]} & $+0.048 \pm 0.032$ & 4/5 \\
& DRL-EntOnly              & 5 & $-1.48$ \tiny{[$-2.65, -0.20$]} & $+0.068 \pm 0.032$ & 4/5 \\
\bottomrule
\end{tabular}
\end{table}

\textbf{Exact-delta Greedy strictly outperforms DRL in slope reduction across all three regions.} The gap (exact-delta Greedy minus DRL-Lagrangian mean) is $-10.68$ pp in Region A, $-9.83$ pp in Region B, and $-3.67$ pp in Region C. The descriptive bootstrap intervals from five seeds do not overlap the exact-delta Greedy value, and the per-seed appendix shows that the gap is not driven by a single outlier. \textbf{However, the slope advantage of exact-delta Greedy comes at a clear contiguity cost}: it returns $-0.55$, $-0.17$, $-0.16$ changes in contiguity across A, B, C. This motivates the scalarised local sweep below rather than a simple slope-only comparison.

\begin{table}[htbp]
\centering
\caption{Contiguity-aware exact-delta local sweep. The table reports three representative points from the seven-weight sweep: the slope-optimal point ($\lambda_{\mathrm{cont}}=0$), the first sampled point with non-negative contiguity change, and the maximum-contiguity point. The final column counts individual deterministic DRL seed outcomes in the same region dominated by at least one sampled local point.}
\label{tab:local-pareto}
\small
\begin{tabular}{llrrrc}
\toprule
Region & Local point & $\lambda_{\mathrm{cont}}$ & Slope $\Delta$\,(\%) & Cont. $\Delta$ & DRL seeds dom. \\
\midrule
\multirow{3}{*}{A} & Slope-optimal & 0 & $-9.66$ & $-0.549$ & \multirow{3}{*}{9/10} \\
& First non-neg. cont. & 1 & $-0.91$ & $+0.250$ & \\
& Maximum contiguity & 8 & $+5.58$ & $+0.405$ & \\
\midrule
\multirow{3}{*}{B} & Slope-optimal & 0 & $-7.64$ & $-0.174$ & \multirow{3}{*}{10/10} \\
& First non-neg. cont. & 0.5 & $-6.19$ & $+0.005$ & \\
& Maximum contiguity & 8 & $+3.70$ & $+0.251$ & \\
\midrule
\multirow{3}{*}{C} & Slope-optimal & 0 & $-5.87$ & $-0.161$ & \multirow{3}{*}{8/10} \\
& First non-neg. cont. & 1 & $-2.83$ & $+0.107$ & \\
& Maximum contiguity & 8 & $+1.54$ & $+0.204$ & \\
\bottomrule
\end{tabular}
\end{table}

The contiguity-aware sweep changes the interpretation of the local-baseline comparison. In every region, at least one sampled local point simultaneously improves slope and preserves or improves contiguity relative to both DRL mean outcomes. The sampled local grid dominates all six DRL mean points and 27 of 30 individual deterministic DRL runs. The three undominated cases are not failures of the test: A EntOnly seed 3 and C seed 1 under both configurations achieve slightly stronger slope reductions than the nearest sampled non-negative-contiguity local point, while retaining positive contiguity. The evidence therefore supports a bounded conclusion: DRL does not Pareto-dominate transparent local ranking in this benchmark, and most of its contiguity benefit is reproducible by an explicit local scalarisation.

\subsection{No-reuse action-mask parity audit}\label{sec:results-mask-parity}

The local sweep in Table~\ref{tab:local-pareto} ranks over the current farmland and forest sets. By contrast, Maskable PPO evaluation uses the environment's converted-parcel mask, so a parcel selected once cannot be selected again later in the same episode. We therefore re-ran the seven-weight local sweep with the same no-reuse constraint to test whether the local advantage depended on a looser feasible set.

\begin{table}[htbp]
\centering
\caption{No-reuse parity audit for the scalarised exact-delta local sweep. The no-reuse version removes both parcels in each committed swap from all later candidate sets, matching the DRL episode-level converted-parcel mask. SE denotes the slope-exact point; NN denotes the sampled non-negative-contiguity point with the strongest slope reduction. Max shift is the largest absolute slope difference between no-reuse and unrestricted local ranking across seven $\lambda_{\mathrm{cont}}$ values.}
\label{tab:masked-local}
\scriptsize
\begin{tabular}{@{}lrrrrrrr@{}}
\toprule
Region & SE $\Delta s$ & SE $\Delta c$ & NN $\lambda$ & NN $\Delta s$ & NN $\Delta c$ & Max shift & Seeds dom. \\
\midrule
A & $-9.66$ & $-0.549$ & 1 & $-0.39$ & $+0.255$ & 2.187 & 6/10 \\
B & $-7.64$ & $-0.174$ & 0.5 & $-6.19$ & $+0.005$ & 0.103 & 10/10 \\
C & $-5.87$ & $-0.161$ & 1 & $-2.83$ & $+0.107$ & 0.000 & 8/10 \\
\bottomrule
\end{tabular}
\end{table}

The no-reuse constraint does not alter the slope-exact result. The $\lambda_{\mathrm{cont}}=0$ outcomes in all three regions are identical to the unrestricted local sweep to the reported precision. The mean-level Pareto conclusion is also unchanged: the no-reuse local grid dominates all six DRL mean outcomes. At the individual-rollout level, the stricter feasible set dominates 24 of 30 deterministic DRL runs, compared with 27 of 30 for the unrestricted grid. Most of the reduction occurs in Region A at high contiguity weights, where the maximum slope shift relative to unrestricted local ranking is 2.187 pp. The audit therefore narrows, rather than weakens, the claim: the main local-vs-DRL conclusion does not rely on an action-mask mismatch, but individual Pareto dominance counts are sensitive to the exact local feasible-set convention.

\subsection{Two-step lookahead audit of local myopia}\label{sec:results-lookahead}

The local sweep above is one-step greedy. A potential objection is that DRL may add value precisely by anticipating downstream swaps that a one-step ranking cannot see. We tested this with a top-5 two-step local lookahead at $\lambda_{\mathrm{cont}}\in\{0,0.5,1\}$, covering the slope-exact setting and the key non-negative-contiguity settings in Table~\ref{tab:local-pareto}. At each paired exchange, the algorithm temporarily evaluates the best next swap after each of the five best immediate candidates before committing the first swap.

\begin{table}[htbp]
\centering
\caption{Two-step lookahead audit for the scalarised exact-delta local ranking. $\Delta$ columns report lookahead minus one-step results. Slope differences are percentage points; negative values mean lookahead improves slope relative to one-step local ranking.}
\label{tab:lookahead}
\scriptsize
\begin{tabular}{lrrrrrr}
\toprule
Region & $\lambda_{\mathrm{cont}}$ & One-step $\Delta s$ & Lookahead $\Delta s$ & $\Delta$ slope pp & One-step $\Delta c$ & Lookahead $\Delta c$ \\
\midrule
A & 0   & $-9.656$ & $-9.656$ & $+0.000$ & $-0.549$ & $-0.549$ \\
A & 0.5 & $-5.927$ & $-5.879$ & $+0.048$ & $-0.014$ & $-0.011$ \\
A & 1   & $-0.910$ & $-0.831$ & $+0.079$ & $+0.250$ & $+0.255$ \\
\midrule
B & 0   & $-7.643$ & $-7.643$ & $+0.000$ & $-0.174$ & $-0.173$ \\
B & 0.5 & $-6.186$ & $-6.191$ & $-0.005$ & $+0.005$ & $+0.007$ \\
B & 1   & $-1.968$ & $-2.122$ & $-0.154$ & $+0.141$ & $+0.138$ \\
\midrule
C & 0   & $-5.871$ & $-5.871$ & $+0.000$ & $-0.161$ & $-0.161$ \\
C & 0.5 & $-5.086$ & $-5.090$ & $-0.004$ & $-0.026$ & $-0.026$ \\
C & 1   & $-2.831$ & $-2.831$ & $+0.000$ & $+0.107$ & $+0.107$ \\
\bottomrule
\end{tabular}
\end{table}

The lookahead audit does not uncover a missed sequential advantage. The slope-exact $\lambda_{\mathrm{cont}}=0$ result is unchanged to three decimals in Regions A and C and changes by only $0.00002$ pp in Region B. Across all nine tested region--weight combinations, the largest absolute slope difference is $0.154$ pp and the largest contiguity difference is $0.0053$. These are small relative to the DRL--local gaps in Table~\ref{tab:main}. The result supports the near-independence interpretation: under these low-budget synthetic-parcel regimes, most attainable benefit is already captured by exact one-step deltas.

\subsection{Budget-ratio sensitivity of local ranking and fixed DRL decoding}\label{sec:results-budget}

The main comparison fixes the intervention budget at $B=100$ paired exchanges. To test whether the local-baseline result is tied to that single budget, we ran two budget-ratio sensitivity checks at $B\in\{25,50,100,200\}$. First, we re-ran the contiguity-aware exact-delta local sweep for all three regions. Second, we freshly decoded the 30 already-trained DRL policies with deterministic paired inference under each budget by setting the episode length to $2B$ raw actions. The second test changes the inference budget only; it is not budget-specific DRL retraining.

\begin{table}[htbp]
\centering
\caption{Budget-ratio sensitivity of local ranking. ``Slope-exact'' is the $\lambda_{\mathrm{cont}}=0$ point. ``Non-neg.'' is the sampled local point with non-negative contiguity change and the strongest slope reduction at that budget.}
\label{tab:budget-sensitivity}
\scriptsize
\begin{tabular}{lrrrrrr}
\toprule
Region & $B$ & Slope-exact $\Delta s$ & Slope-exact $\Delta c$ & Non-neg. $\lambda$ & Non-neg. $\Delta s$ & Non-neg. $\Delta c$ \\
\midrule
A & 25  & $-5.05$ & $-0.170$ & 0.5 & $-3.87$ & $+0.012$ \\
A & 50  & $-7.33$ & $-0.329$ & 0.5 & $-5.27$ & $+0.002$ \\
A & 100 & $-9.66$ & $-0.549$ & 1.0 & $-0.91$ & $+0.250$ \\
A & 200 & $-9.89$ & $-0.653$ & 1.0 & $-0.91$ & $+0.250$ \\
\midrule
B & 25  & $-3.24$ & $-0.048$ & 0.5 & $-2.70$ & $+0.019$ \\
B & 50  & $-5.14$ & $-0.083$ & 0.5 & $-4.23$ & $+0.023$ \\
B & 100 & $-7.64$ & $-0.174$ & 0.5 & $-6.19$ & $+0.005$ \\
B & 200 & $-9.57$ & $-0.287$ & 1.0 & $-1.75$ & $+0.148$ \\
\midrule
C & 25  & $-2.43$ & $-0.052$ & 0.5 & $-2.11$ & $+0.006$ \\
C & 50  & $-3.83$ & $-0.088$ & 1.0 & $-2.41$ & $+0.054$ \\
C & 100 & $-5.87$ & $-0.161$ & 1.0 & $-2.83$ & $+0.107$ \\
C & 200 & $-8.30$ & $-0.287$ & 1.0 & $-3.25$ & $+0.172$ \\
\bottomrule
\end{tabular}
\end{table}

Two local-frontier patterns are relevant for decision support. First, the slope-exact local point improves slope more strongly as $B$ increases in all three regions, but the contiguity loss also deepens. Second, a sampled local point with non-negative contiguity exists in every tested region-budget pair and still gives negative slope change in all 12 cases. The local-frontier result is therefore not an artefact of choosing $B=100$.

\begin{table}[htbp]
\centering
\caption{Fixed-policy DRL decoding at alternative paired-exchange budgets compared with the slope-exact local baseline. $\Delta s$ values are percentage changes in area-weighted farmland slope. ``Best DRL'' is the most negative slope outcome among the 10 fixed policies in that region-budget pair.}
\label{tab:drl-budget-decode}
\scriptsize
\begin{tabular}{lrrrrr}
\toprule
Region & $B$ & DRL mean $\Delta s$ & Best DRL $\Delta s$ & Best DRL $\Delta c$ & Slope-exact local $\Delta s$ \\
\midrule
A & 25  & $+0.35$ & $-0.53$ & $+0.110$ & $-5.05$ \\
A & 50  & $+0.17$ & $-1.08$ & $+0.168$ & $-7.33$ \\
A & 100 & $+0.83$ & $-0.97$ & $+0.039$ & $-9.66$ \\
A & 200 & $+4.08$ & $+1.37$ & $-0.145$ & $-9.89$ \\
\midrule
B & 25  & $+0.24$ & $-0.05$ & $-0.044$ & $-3.24$ \\
B & 50  & $+0.73$ & $+0.36$ & $-0.080$ & $-5.14$ \\
B & 100 & $+2.19$ & $+1.62$ & $-0.159$ & $-7.64$ \\
B & 200 & $+6.06$ & $+4.97$ & $-0.302$ & $-9.57$ \\
\midrule
C & 25  & $-0.63$ & $-1.05$ & $+0.025$ & $-2.43$ \\
C & 50  & $-1.22$ & $-2.25$ & $+0.020$ & $-3.83$ \\
C & 100 & $-1.84$ & $-3.33$ & $+0.037$ & $-5.87$ \\
C & 200 & $-2.23$ & $-4.45$ & $+0.047$ & $-8.30$ \\
\bottomrule
\end{tabular}
\end{table}

The fixed-policy decoding test strengthens, rather than weakens, the local-baseline conclusion. The slope-exact local baseline outperforms the best DRL rollout in all 12 region-budget pairs. Regions A and B show clear budget fragility under deterministic decoding: increasing the budget worsens mean slope outcomes, especially at $B=200$. Region C improves with larger budgets, but even its best individual DRL rollout remains less effective than slope-exact local ranking. This result closes the fixed-policy decoding concern, but not the retraining question: policies trained specifically for a different budget could still behave differently.

Because the policy observation includes a normalised progress feature, we also decoded $B=25$ and $B=50$ as prefixes of the training-time 100-pair horizon rather than as native shorter episodes. This protocol check preserved the conclusion: slope-exact local ranking still beat the best DRL rollout in all six smaller-budget region-budget pairs, while DRL mean slope outcomes changed by at most 0.052 percentage points relative to native-horizon decoding.

Figure~\ref{fig:cross_region} visualises the comparison.

\begin{figure}[htbp]
\centering
\includegraphics[width=0.85\textwidth]{figures/fig1_cross_region.png}
\caption{Cross-region comparison of primary methods under area-weighted paired-inference evaluation. The figure separates slope-only local ranking from exact-delta local ranking. Exact-delta local ranking gives the strongest slope reduction in all three regions, while DRL configurations regress slope in Regions A and B and improve it modestly in Region C. Deterministic local baselines are shown without error bars; DRL error bars show across-seed standard deviations. The contiguity cost of exact-delta ranking is evaluated separately in Table~\ref{tab:local-pareto} and Figure~\ref{fig:pareto}.}
\label{fig:cross_region}
\end{figure}

\subsection{Per-seed bimodal convergence}\label{sec:results-bimodal}

Figure~\ref{fig:bimodal} disaggregates DRL outcomes per seed. Region A exhibits a strikingly bimodal pattern: seeds 0--3 cluster near zero (between $-1.0\%$ and $+0.4\%$ slope change), while seed 4 produces $+5.17\%$ to $+5.25\%$---a clear failure mode---in both Lagrangian and EntOnly configurations. The fact that the same RNG seed (4) fails under both reward-weighting strategies establishes that bimodal convergence is a property of the optimisation landscape itself, not of the reward shaping. Region B shows uniform failure (all 10 policies produce positive slope), and Region C shows mostly negative slope with one positive outlier (EntOnly seed 3, $+0.88\%$).

\begin{figure}[htbp]
\centering
\includegraphics[width=\textwidth]{figures/fig2_per_seed_bimodal.png}
\caption{Per-seed paired-inference outcomes. Region A is bimodal, with four seeds clustering near zero slope change (approximately $-1.0\%$ to $+0.4\%$) and seed 4 at $+5\%$. Region B shows uniform deterministic failure: every seed in both configurations regresses slope. Region C contains mostly modestly improving seeds and one positive outlier. Dotted green and dashed red reference lines show slope-only and exact-delta local ranking, respectively.}
\label{fig:bimodal}
\end{figure}

\subsection{Train--evaluation distribution mismatch in region B}\label{sec:results-mismatch}

The most striking finding is Region B's uniform DRL failure under deterministic paired-inference (all 10 trained policies have slope changes between $+1.57\%$ and $+2.36\%$), despite training-time stochastic rollouts reporting consistent improvement of the area-weighted slope objective. Both training and evaluation enforce paired action masking and farmland-count conservation; the difference is in action-distribution sharpness. During training, the agent samples actions stochastically from the masked softmax distribution---generating diverse rollouts that explore many low-probability parcels. During paired-inference evaluation, the agent commits to the $\arg\max$ of the same masked distribution, yielding a single deterministic trajectory. The mismatch hypothesis posits that the policy converges to a softmax distribution that produces favourable rollouts \emph{in expectation} (when stochastically sampled with sufficient entropy) but yields a single failing trajectory under deterministic argmax. Section~\ref{sec:results-stochastic} tests this hypothesis directly.

Three structural features may amplify this asymmetry in Region B: the highest farmland-to-forest ratio (2.98 vs.\ 2.40 in A and 2.87 in C), constraining forest-pool diversity; the lowest CV(area) of the three synthetic regions; and the lowest mean Queen degree (4.47 vs.\ 5.28/5.61). Each of these tightens the action-decision coupling in ways that compress the stochastic exploration benefit, though disentangling them within three regions is not possible.

\subsection{Stochastic vs deterministic paired-inference}\label{sec:results-stochastic}

If train-evaluation distribution mismatch alone drove Region B's failure, evaluating the same trained policies under stochastic-sampling paired-inference (matching training-time action selection) should recover negative slope outcomes. We test this directly on all 10 Region B policies trained under the area-weighted environment. Each policy is evaluated twice: (a) the deterministic argmax protocol used in Table~\ref{tab:main}, and (b) stochastic sampling from the masked softmax, repeated for 5 trials per policy and averaged.

\begin{table}[htbp]
\centering
\caption{Stochastic-sampling vs deterministic-argmax paired-inference on Region B (area-weighted env). Stochastic inference improves mean slope by approximately 1.3 percentage points relative to argmax, but does not produce negative slope: it remains far from the slope-only local baseline at $-4.30\%$, the slope-exact local baseline at $-7.64\%$, and the non-negative-contiguity local point at $-6.19\%$. Mismatch is a contributing factor, not the sole driver.}
\label{tab:stochastic}
\begin{tabular}{lllrr}
\toprule
Cfg & Inference & $n$ & Slope $\Delta$\,(\%) (mean $\pm$ std) & Range \\
\midrule
\multirow{2}{*}{Lagrangian} & Argmax (deterministic) & 5 & $+2.22 \pm 0.28$ & $[+1.68, +2.45]$ \\
                            & Stochastic (5 trials/policy) & 5 & $+0.86 \pm 0.05$ & $[+0.76, +0.92]$ \\
\midrule
\multirow{2}{*}{EntOnly}    & Argmax (deterministic) & 5 & $+2.17 \pm 0.28$ & $[+1.62, +2.40]$ \\
                            & Stochastic (5 trials/policy) & 5 & $+1.09 \pm 0.57$ & $[+0.68, +2.23]$ \\
\bottomrule
\end{tabular}
\end{table}

The results are informative. Switching from argmax to stochastic inference reduces the mean slope change from $+2.22\%$ to $+0.86\%$ for Lagrangian (a $1.36$ pp improvement) and from $+2.17\%$ to $+1.09\%$ for EntOnly. However, the stochastic outcome remains positive in every policy: no policy under either configuration achieves negative slope. Slope-only Greedy's $-4.30\%$ slope reduction, slope-exact Greedy's $-7.64\%$ reduction, and the non-negative-contiguity local point's $-6.19\%$ reduction are all substantially better than the best stochastic-inference DRL outcome.

Two findings follow. First, train-evaluation distribution mismatch is a real and quantifiable contributor to the DRL-vs-local-baseline gap; about $1.3$ percentage points of the gap to slope-only Greedy can be attributed to it. Second, mismatch is \emph{not} the sole driver: even when matched, DRL fails to produce negative slope, leaving a large residual gap to slope-only, slope-exact, and contiguity-aware local ranking. The most plausible candidates for the residual are the weak inter-step coupling on irregular parcels (Section~\ref{sec:results-overlap}), and the area-weighted reward landscape's relatively flat structure beyond the slope-greedy ordering. One EntOnly policy (seed 4) shows essentially identical outcomes under argmax and stochastic inference ($+2.24\%$ vs $+2.23\%$), suggesting it has collapsed to a near-deterministic distribution---another instance of bimodal failure mode (Section~\ref{sec:results-bimodal}).

\subsection{Inference trajectories}\label{sec:results-traj}

Figure~\ref{fig:traj} plots the step-by-step slope evolution during paired-inference. Region A's converged seeds (0--3) show smooth descent below the initial slope; seed 4 shows the canonical failure pattern of monotone ascent. Region B shows 10 trajectories starting at the initial slope and rising---no seed produces a downward trend. Region C trajectories descend modestly but with reduced rate compared to region A.

\begin{figure}[htbp]
\centering
\includegraphics[width=\textwidth]{figures/fig3_trajectories.png}
\caption{Step-by-step area-weighted slope evolution during deterministic 200-step paired-inference. Faint traces show individual seeds; bold traces highlight seed 0 of each configuration. Region A shows both convergent and divergent trajectories, Region B shows monotone divergence across all policies, and Region C shows modest descent with one divergent trajectory.}
\label{fig:traj}
\end{figure}

\subsection{Slope--contiguity tradeoff}\label{sec:results-pareto}

Figure~\ref{fig:pareto} positions the primary methods and the seven-weight local sweep in the slope--contiguity plane. The ideal location is the upper-left quadrant (lower slope, higher contiguity). The slope-only local point lies in the upper-left in all three regions; exact-delta local ranking lies in the lower-left (strongest slope but degraded contiguity); GA stays near the origin; and DRL spreads broadly. The scalarised local sweep connects these local extremes and shows that positive-contiguity slope reductions can be recovered without a learned policy. In Region B, DRL points fall in the lower-right quadrant (worse on both axes), confirming the failure mode of Section~\ref{sec:results-mismatch}.

\begin{figure}[htbp]
\centering
\includegraphics[width=\textwidth]{figures/fig4_tradeoff.png}
\caption{Slope--contiguity Pareto plane for primary comparators and the scalarised exact-delta local sweep. Per-seed DRL outcomes are plotted individually; deterministic local baselines are single points. The local sweep fills the gap between slope-only and slope-exact local ranking and contains positive-contiguity slope-reducing points in all three regions. Region B's DRL points cluster in the lower-right quadrant, indicating degradation on both objectives.}
\label{fig:pareto}
\end{figure}

\subsection{Computational cost}

Per-run wall time on the same hardware: Greedy $<\!1$\,s; Random $<\!1$\,s; GA 15--35\,s (5 seeds, $\sim$5{,}000 fitness evaluations); DRL 1.5--3.5\,h on NVIDIA A100. The 21-run contiguity-aware local sweep completed in approximately 4 minutes in the reproducibility implementation, including analytic-delta validation checks. The DRL training cost across 30 models was approximately 75 GPU-hours; baseline computation remains several orders of magnitude cheaper than DRL training.

\subsection{Empirical neighbour-overlap rate}\label{sec:results-overlap}

The near-independence hypothesis (Section~\ref{sec:disc-near-indep}) predicts that consecutive paired exchanges affect approximately disjoint neighbourhoods on irregular cadastral parcels. We test this by replaying every paired-inference rollout (30 episodes $\times$ 100 paired exchanges) and recording, for each consecutive pair of paired exchanges $(p_k, p_{k+1})$, whether $p_{k+1}$ touches any Queen neighbour of either parcel involved in $p_k$. The null distribution is estimated by Monte Carlo: 1{,}000 random rollouts per region under the same enforced paired-action-mask and farmland-count-conservation constraints, using uniform random selection from the masked pool at each step.

\begin{table}[htbp]
\centering
\caption{Empirical neighbour-overlap rate (DRL paired-inference) versus a Monte Carlo null that respects action masks, land-use pools, and the paired-budget constraint. Empirical means are across all 10 DRL policies per region (5 seeds $\times$ 2 configurations); MC null is from 1{,}000 random rollouts. The empirical rate exceeds the MC null 97.5 percentile by orders of magnitude in every region.}
\label{tab:overlap}
\begin{tabular}{lrrrr}
\toprule
Region & MC null mean & MC null 95\% CI & DRL empirical mean & DRL/null ratio \\
\midrule
A (1{,}601 swap) & 0.93\% & $[0.00, 3.03]$ & $43.9 \pm 5.1$\% & $\sim 47\times$ \\
B (2{,}669 swap) & 0.55\% & $[0.00, 2.02]$ & $37.9 \pm 4.2$\% & $\sim 69\times$ \\
C (6{,}787 swap) & 0.28\% & $[0.00, 2.02]$ & $17.8 \pm 3.7$\% & $\sim 64\times$ \\
\bottomrule
\end{tabular}
\end{table}

Two findings emerge. First, DRL's empirical overlap rate exceeds the Monte Carlo null mean by 47--69$\times$ and lies well outside the null 97.5 percentile in every region: the policy has reliably learned a spatially clustered strategy, contrary to a strict null-hypothesis interpretation of near-independence. Second, the absolute overlap rate remains modest---17.8\% in Region C, 37.9--43.9\% in Regions A and B---meaning that more than half of consecutive paired exchanges still affect disjoint neighbourhoods. Even where DRL has learned to cluster, the resulting sequential dependence is insufficient to overcome local-ranking strategies under area-weighted aggregate objectives: aggregate-metric changes are still dominated by intrinsic parcel properties (Section~\ref{sec:disc-near-indep}). The combination of strong learned clustering and yet-still-insufficient inter-step coupling explains why DRL produces above-random but below-Greedy outcomes.




% ============================================================================
\section{Discussion}\label{sec:disc}
% ============================================================================

\subsection{Near-independence hypothesis, refined}\label{sec:disc-near-indep}

The principal observation is that exact one-step local ranking captures more of the tested decision-support objective than the DRL policies do. The slope-only exact-delta baseline already outperforms DRL on area-weighted slope, and the contiguity-aware local sweep shows that this is not merely an unfair comparison between a single-objective heuristic and a multi-objective policy: sampled local points dominate all six DRL mean outcomes and 27 of 30 individual deterministic DRL runs. The no-reuse parity audit in Section~\ref{sec:results-mask-parity} shows that the mean-level conclusion does not rely on a looser local feasible set, and the two-step lookahead audit in Section~\ref{sec:results-lookahead} further shows that the result is not an artefact of purely myopic one-step selection. The \emph{swap near-independence hypothesis} posits that individual swap decisions become weakly coupled when parcel areas are heterogeneous, adjacency is sparse, and the exchange budget is small relative to the candidate pool. The empirical analyses in Sections~\ref{sec:results-overlap}, \ref{sec:results-stochastic}, \ref{sec:results-main}, \ref{sec:results-pareto}, \ref{sec:results-mask-parity}, and~\ref{sec:results-lookahead} qualify this hypothesis in several ways.

\textbf{Two-step lookahead adds little.} If sequential coupling were strong, a top-5 two-step lookahead should have shifted the local frontier materially. It did not: the largest slope change relative to one-step local ranking is $0.154$ pp, and the slope-exact setting is effectively unchanged. This does not prove global optimality. It does show that the tested local frontier is not a fragile consequence of using a purely myopic ranking.

\textbf{Action-mask parity does not remove the local advantage.} Imposing the same no-reuse converted-parcel constraint used by DRL leaves the slope-exact local results unchanged in all three regions. It also preserves dominance over all six DRL mean outcomes, although individual-rollout dominance falls from 27/30 to 24/30. This distinction matters: the result is robust at the level of the paper's mean comparison, but individual Pareto counts should be interpreted as convention-dependent diagnostics rather than universal rankings.

\textbf{Spatial sparsity is real but partial.} With $B = 100$ exchanges among thousands of parcels and average Queen degree $\sim$5, the uniform-random expectation for two consecutive swaps to share a neighbour is approximately $2 \times 5.61 / 6{,}787 \approx 0.17\%$ in Region C. The empirical neighbour-overlap rates (Section~\ref{sec:results-overlap}) exceed this random baseline by two orders of magnitude (18--44\%): DRL has learned a spatially clustered strategy. Yet the absolute overlap remains below 50\%, leaving more than half of consecutive exchanges in disjoint neighbourhoods. The inter-step dependence DRL achieves is real but bounded.

\textbf{Train-evaluation distribution mismatch is a partial driver.} The stochastic-vs-deterministic comparison in Section~\ref{sec:results-stochastic} shows that switching from argmax to stochastic-sampling inference improves Region B's mean slope outcome by approximately 1.3 percentage points, closing about 20\% of the slope-only-Greedy gap. Mismatch is therefore a real and quantifiable contributor, but it does not by itself produce negative slope: even the matched stochastic outcome is positive ($+0.86\%$ to $+1.09\%$) while slope-only Greedy reaches $-4.30\%$. The remaining $\sim$80\% reflects deeper limits.

\textbf{The full slope gap to exact-delta local ranking is substantially larger than the stochastic-inference recovery covers.} On Region B, slope-exact Greedy reaches $-7.64\%$ slope reduction, and the contiguity-aware local point at $\lambda_{\mathrm{cont}}=0.5$ reaches $-6.19\%$ while keeping contiguity slightly positive. The best stochastic-inference DRL outcome is about $+0.86\%$. Closing this gap is not within reach of inference-time decoding alone. It would require a policy class, training objective, or pair-scoring architecture that represents the marginal area-weighted and contiguity objectives more directly than the tested per-parcel scorer.

\textbf{The contiguity tradeoff is not uniquely solved by DRL.} The original slope-exact Greedy point sacrifices contiguity, but explicit local scalarisation recovers positive contiguity in all three regions. This does not mean every DRL run is dominated: three individual deterministic DRL runs remain outside the unrestricted sampled local grid, and six remain outside the stricter no-reuse grid. It does mean that DRL's apparent multi-objective advantage is not sufficient evidence of a learned sequential advantage unless it is compared against a transparent multi-objective local sweep.

\subsection{Why region B fails uniquely under paired-inference}\label{sec:disc-B}

Region B's failure of all 10 trained policies under paired-inference---despite consistent training-rollout improvement---reflects a structural train--evaluation mismatch unique to the most-constrained intermediate scale. Both training and evaluation enforce paired action masking and farmland-count conservation; the difference is in action-distribution sharpness (Section~\ref{sec:results-mismatch}). The stochastic-vs-deterministic comparison (Section~\ref{sec:results-stochastic}) shows that this mismatch closes about 1.3 percentage points of the gap to slope-only Greedy, so it is a real but partial driver. Two structural features of Region B may amplify the residual gap:

\paragraph{Forest-pool scarcity.} Region B has the highest farmland-to-forest ratio (2.98). With $B = 100$ paired exchanges and only $\sim$671 forest parcels available, the forest pool is small enough that any single argmax-selected forest parcel determines a disproportionately large share of the trajectory. Stochastic sampling smooths this concentration but does not eliminate it.

\paragraph{Low CV and degree.} Region B's CV(area) and mean Queen degree (4.47) are the lowest of the three regions. Lower degree concentrates contiguity rewards on a few neighbours, sharpening the optimal action distribution under training reward but increasing sensitivity to argmax decoding at evaluation time.

This result exemplifies a broader concern in DRL applications under safety or feasibility constraints: a policy that performs well in expectation under stochastic sampling does not automatically perform well under deterministic argmax---particularly when forest-pool size and adjacency degree push the optimal distribution to be sharply peaked. Inference-time decoding strategies that retain entropy (temperature sampling, top-$k$, beam search) may partially close this residual gap.

\subsection{Comparison with grid proof-of-concept}\label{sec:disc-vs-grid}

Published DRL studies for routing, allocation, and grid-like spatial planning \citep{bello2017neural,nazari2018reinforcement,kool2019attention,zheng2023spatial} report strong results under regular or highly structured decision spaces. Why does the same family of policy-gradient pipelines struggle on synthetic parcels of comparable size?

The structural difference is plausible and measurable. Grid cells have uniform area (CV = 0) and dense, symmetric neighbourhoods (8 neighbours each in Queen contiguity). Both properties weaken the near-independence assumption: equal-area swaps create more uniformly coupled aggregate effects, and dense adjacency creates stronger contiguity interactions between swap decisions. In a $20{\times}20$ grid with $B = 100$ swaps among 200 cells, the exchange-budget ratio $B / N_{\text{swap}} = 0.5$ is two orders of magnitude higher than in the present synthetic regions ($B / N_{\text{swap}} = 0.015$--$0.062$), creating denser inter-swap interactions.

Grid abstractions therefore operate in a regime where DRL's sequential-decision advantage is more likely to be expressed. The transition from grids to irregular parcels is not merely a scaling challenge. It changes the optimisation landscape through area heterogeneity, sparse adjacency, and low intervention ratios. This interpretation remains bounded to the tested Maskable PPO and scorer-MLP family.

\subsection{Reproducibility through synthetic data}\label{sec:disc-repro}

Beyond the methodological findings, this study contributes a reproducible synthetic pipeline that addresses the access restrictions hampering Chinese cadastral GIS research. The pipeline (released in an anonymised review repository) generates heuristically matched synthetic regions from Copernicus DEM and Impact Observatory land-cover alone---no restricted cadastral boundaries are consumed. The 30 trained models, paired-inference evaluations, and baseline runs are released as artefacts.

A direct benefit is that researchers without TNLS access can now benchmark land-use optimisation algorithms against documented synthetic-cadastral characteristics. A second benefit is that structural parameters---CV(area), avg-degree, farm:forest ratio---can be parametrically varied within the pipeline, supporting hypothesis-testing experiments that real-data case studies cannot.

\subsection{Implications for practice and future work}\label{sec:disc-implications}

Three practical implications follow:

\begin{table}[htbp]
\centering
\caption{Diagnostic workflow for deciding whether parcel-level DRL is justified in land-use decision support. The workflow is intended as a screening procedure before GPU-intensive training, not as a replacement for local planning judgement.}
\label{tab:diagnostic-workflow}
\small
\begin{tabular}{p{0.19\textwidth}p{0.36\textwidth}p{0.36\textwidth}}
\toprule
Stage & Computation & Decision-support use \\
\midrule
1. Characterise the planning area & Compute CV(area), average Queen degree, farm:forest ratio, candidate-pool size, and exchange-budget ratio $B/N_{\text{swap}}$. & Identify whether sparse adjacency, heterogeneous parcel areas, or low intervention ratios make near-independent local swaps plausible. \\
2. Run transparent local diagnostics & Evaluate slope-only ranking, exact-delta slope ranking, and a scalarised slope--contiguity local sweep. & Establish the attainable local frontier and the cost of preserving contiguity before training any learned policy. \\
3. Compare with planning thresholds & Inspect whether a local point satisfies the acceptable slope--contiguity tradeoff for the planning task. & Stop at local ranking if the transparent frontier already meets the management need. \\
4. Escalate selectively & Train DRL or stronger metaheuristics only when the local frontier fails or when diagnostics suggest strong inter-swap coupling. & Reserve costly learned models for cases where sequential dependence is likely to add value. \\
5. Report audit artefacts & Release seeds, per-run outcomes, budget checks, and code needed to reproduce the comparison. & Make the algorithm choice transparent to reviewers, planners, and downstream users. \\
\bottomrule
\end{tabular}
\end{table}

\textbf{Multi-objective local-ranking baselines should be mandatory diagnostics before deploying DRL.} On the synthetic-parcel structures tested here, slope-exact Greedy produces stronger slope reductions than any DRL configuration while running in seconds rather than GPU-hours. More importantly, a seven-weight contiguity-aware exact-delta sweep recovers positive contiguity in all three regions and dominates all DRL mean outcomes. Table~\ref{tab:diagnostic-workflow} translates this finding into a screening workflow for CEUS-style decision support. We do not claim that DRL is dominated in every cadastral setting, only that it failed to Pareto-dominate transparent local baselines in the regimes tested.

\textbf{Algorithm choice should follow problem structure.} DRL is most appropriate where sequential decisions exhibit genuine inter-step dependencies. In irregular parcel systems with low $B / N_{\text{swap}}$ ratios, this condition should not be assumed. Practitioners should measure CV(area), adjacency degree, and exchange-budget ratio before deploying DRL.

\textbf{Block-level reformulation is a testable next direction.} If DRL is to provide value on irregular cadastral-style data, the problem may need to be reformulated at higher levels of spatial abstraction where stronger dependencies exist. Operating at the level of contiguous farmland blocks ($\geq$6.67~ha), aligned with China's \emph{tiankuai} administrative unit, would reduce the action space from thousands of parcels to tens or hundreds of blocks while creating resource-allocation tradeoffs across competing blocks. This is a hypothesis for subsequent work, not a conclusion established by the present experiments.

\subsection{Limitations}\label{sec:disc-limits}

\textbf{CV(area) below target.} The SLIC tessellation produces parcels with CV(area) of 0.57--0.85, below the published target of $\sim$2.0 for real cadastral data. SLIC's compactness objective explicitly encourages superpixel size uniformity \citep{achanta2012slic}, suppressing CV(area) below what real cadastral histories produce. The post-hoc trim-to-target step (Section~\ref{sec:synth}) restores parcel-count targets but does not recover area heterogeneity. We do not have evidence about how the DRL-local-baseline comparison scales under higher CV(area); a multi-seed phase-transition experiment under the area-weighted environment is left to future work.

\textbf{Single geographic envelope.} All three regions are sampled from a $\sim$50$\times$50~km extent in transitional Sichuan Basin terrain. Cadastral structures in alluvial plains, coastal regions, or northern China may differ in ways that affect the near-independence regime. The pipeline supports re-sampling other extents.

\textbf{Convergence threshold heuristic.} The 1\% slope-reduction convergence threshold is a heuristic; alternative thresholds shift convergence rates marginally but do not affect the headline finding that exact-delta Greedy achieves stronger slope reduction than DRL means in all three tested regions.

\textbf{Budget-specific DRL retraining not tested.} Section~\ref{sec:results-budget} shows that the local slope--contiguity frontier remains informative for $B\in\{25,50,100,200\}$ and that fixed-policy DRL decoding does not close the slope-exact local gap under these budgets. However, the policies were originally trained under the $B=100$ setting. We therefore do not establish how the DRL-local-baseline gap changes if policies are retrained with budget-specific horizons, reward schedules, or curricula. The near-independence interpretation predicts stronger sequential coupling at larger budgets, making matched DRL retraining across budgets a high-priority future sensitivity analysis.

\textbf{Baseline coverage.} The deterministic local-ranking baselines are the principal point of comparison. Genetic Algorithm is included as a conventional stochastic baseline drawn from the land-use optimisation literature, not as an exhaustively tuned optimiser; with population 50 and 100 generations ($\sim$5{,}000 evaluations), the GA budget is modest for cadastral search spaces of $3{,}000$+ parcels. The contiguity-aware local sweep tests seven scalarisation weights, not a continuous Pareto frontier or a fully tuned local-search/metaheuristic family. The no-reuse parity audit matches the DRL converted-parcel mask and confirms the mean-level result, but it remains a deterministic local diagnostic. The two-step lookahead audit reduces the concern that one-step ranking is trivially myopic, but it is still not a global optimiser. A more extensively tuned metaheuristic, such as GA with $10^5$+ evaluations, NSGA-II, simulated annealing, or tabu search, is still needed before claiming the best attainable classical frontier. The present study therefore makes a narrower claim: the sampled Maskable-PPO policies fail to beat transparent exact-delta local diagnostics under the tested regimes, not every possible heuristic.

\textbf{Algorithm coverage and architecture.} We test only Maskable PPO with two reward-weighting variants and a parcel-permutation-invariant MLP policy that consumes hand-engineered per-parcel and global features. Adjacency information is supplied through summary statistics (proportion of same-type neighbours, neighbour count, mean neighbour slope) rather than as a learnable graph encoding. Alternative architectures may capture inter-parcel dependencies that the present scorer cannot, particularly: (i) graph-neural-network policies that propagate features along the Queen adjacency graph; (ii) pair-scoring networks that directly score (farmland, forest) swap pairs rather than parcels in isolation; (iii) attention/pointer-network policies in the style of \citet{kool2019attention} that learn pairwise compatibility. Each could plausibly close part of the gap we report, particularly the train-evaluation mismatch that emerges when sharp action distributions are sampled stochastically and argmax-decoded. Alternative DRL algorithms (SAC, TD3) and inference-time decoding strategies (temperature-controlled sampling, top-$k$, beam search) are similarly untested. We frame our findings as evidence about this specific Maskable-PPO + scorer-MLP family, not as a general claim about all parcel-level DRL.

\textbf{Cadastral generality.} The synthetic-data pipeline parametrises area heterogeneity, average degree, and farm-to-forest ratio independently, but does not model parcel-shape extreme irregularity or fragmentation patterns that are characteristic of some real cadastres. Stronger calibration evidence---area distribution, shape-index distribution, Queen-degree distribution, slope-area correlation, and parcel-pattern maps---is collected in the supplementary release for transparency.

% ============================================================================
\section{Conclusions}\label{sec:concl}
% ============================================================================

This study evaluated whether Maskable PPO provides sufficient additional value for cadastral-style land-use decision support to justify its training cost relative to transparent local baselines. Using a fully reproducible synthetic-parcel pipeline heuristically matched to published aggregate structural ranges for mountainous cadastral settings, 30 Maskable PPO policies were trained across three regions and two reward-weighting strategies, evaluated under deterministic paired-inference with an area-weighted slope objective, and supplemented by a contiguity-aware local sweep, a no-reuse action-mask parity audit, a two-step local lookahead audit, an empirical neighbour-overlap analysis, and a stochastic-vs-deterministic inference comparison. The findings are:

\begin{enumerate}
\item \textbf{Exact-delta Greedy outperforms Maskable PPO on slope across all three regions.} Descriptive bootstrap intervals and raw per-seed outcomes both place the DRL results well away from the exact-delta Greedy value. The slope gap (exact-delta Greedy minus DRL-Lagrangian mean) is $-10.68$ pp in Region A, $-9.83$ pp in Region B, and $-3.67$ pp in Region C.

\item \textbf{The multi-objective tradeoff is largely captured by local scalarisation.} Exact-delta Greedy substantially degrades contiguity (changes of $-0.55$, $-0.17$, $-0.16$ across A/B/C), but the contiguity-aware local sweep finds positive-contiguity points in all three regions. The sampled local grid dominates all six DRL mean outcomes and 27 of 30 individual deterministic DRL runs, while three individual DRL runs remain undominated.

\item \textbf{The local advantage is not explained by an action-mask mismatch.} A no-reuse version of the seven-weight local sweep, matching the DRL converted-parcel mask, leaves the slope-exact outcomes unchanged and still dominates all six DRL mean outcomes. Under this stricter convention, individual-run dominance is 24 of 30 rather than 27 of 30.

\item \textbf{Two-step local lookahead does not materially change the local frontier.} A top-5 two-step lookahead at $\lambda_{\mathrm{cont}}\in\{0,0.5,1\}$ changes slope outcomes by at most $0.154$ pp and contiguity by at most $0.0053$ relative to one-step local ranking. The slope-exact result is effectively unchanged.

\item \textbf{The local frontier and fixed-policy DRL gap are not specific to the 100-pair budget.} Across $B\in\{25,50,100,200\}$, the slope-exact local point improves slope more strongly as budget increases but loses more contiguity, while a sampled non-negative-contiguity local point remains available in every region-budget pair. Fresh deterministic decoding of the 30 trained DRL policies at the same budgets leaves slope-exact local ranking ahead of the best DRL rollout in all 12 region-budget pairs.

\item \textbf{Region B exhibits a clear train-evaluation distribution mismatch}: stochastic-sampling paired-inference (matching training-time action selection) reduces the slope outcome by approximately 1.3 percentage points relative to deterministic argmax, but does not close the full gap to slope-only, slope-exact, or contiguity-aware local ranking. Mismatch is a real and quantifiable contributor (about 20\% of the slope-only-Greedy gap), but not the sole driver.

\item \textbf{Bimodal seed convergence in Region A is reward-weighting--invariant}: the same RNG seed fails under both Lagrangian and fixed-weight configurations, establishing that the bimodal pattern reflects the optimisation landscape itself rather than reward shaping.

\item \textbf{DRL has learned non-trivial spatial clustering, but it is not enough.} Empirical neighbour-overlap rates (18--44\%) exceed the random baseline by two orders of magnitude across all 30 paired-inference rollouts, yet remain below 50\%. The inter-step coupling DRL achieves is real but bounded.
\end{enumerate}

These results identify conditions under which transparent local ranking remains difficult to beat: low exchange-budget ratios, sparse adjacency, and aggregate objectives whose exact one-step deltas can be evaluated cheaply. DRL may still provide value under architectures or abstractions that encode stronger spatial dependencies, but that claim requires separate evidence. The immediate decision-support implication is narrower and stronger: slope-only, slope-exact, and contiguity-aware exact-delta local baselines should be run before any GPU-intensive DRL training pipeline is deployed for cadastral-style layout optimisation.

\section*{Data and code availability}

The synthetic data generation pipeline, all 30 trained models, paired-inference evaluation outputs, baseline results, contiguity-aware local-sweep outputs, no-reuse local parity-audit outputs, two-step local-lookahead outputs, local budget-sensitivity outputs, fixed-policy DRL budget-decoding outputs, prefix-horizon decoding audit outputs, figure-generation audit outputs, and figures are available in an anonymised repository for review at \url{https://anonymous.4open.science/r/cadastral-drl-opt}; should the anonymous mirror be unavailable during review, reviewers may request access through the journal's editorial system. A public non-anonymous archive should replace the review mirror after acceptance. The synthetic data pipeline runs on CPU in approximately 25~minutes; the 21-run contiguity-aware local sweep completes in approximately 4 minutes, the 42-run no-reuse parity audit in approximately 4.6 minutes, the 9-run two-step lookahead audit in approximately 8.5 minutes, the 84-run local budget-sensitivity sweep in approximately 15 minutes, the 120-run fixed-policy DRL budget-decoding sweep in approximately 2 minutes, and the 60-run prefix-horizon audit in under 1 minute in the reproducibility implementation; full DRL training requires approximately 75 GPU-hours on an NVIDIA A100.

\section*{Declaration of generative AI use}

The authors used a large-language-model assistant (Claude, Anthropic) for editorial polishing of figure captions and prose clarity, and for boilerplate-code refactoring of the synthetic-data pipeline. All technical content, experimental design, statistical analyses, and scientific interpretations are the authors' own; the assistant did not contribute novel results or interpretations.

\section*{Disclosure statement}

The authors report no conflict of interest.

\section*{Funding}

Funding information is provided on the separate title page in accordance with double-anonymous review requirements.

\clearpage

% ============================================================================
\appendix
\section{Per-seed paired-inference detail}\label{app:per-seed}
% ============================================================================

Table~\ref{tab:per-seed} reports the paired-inference outcome of every trained policy individually. Initial slope $\bar{s}$ and contiguity $\bar{c}$ are environment constants per region; final values reflect the deterministic policy rollout. All 30 policies satisfy farmland-count conservation ($\Delta n = 0$) and complete the full 100-pair budget.

\begin{table}[!ht]
\centering
\caption{Per-seed paired-inference outcomes for all 30 trained DRL policies under the area-weighted environment. Each row is one trained model; metrics are from a single deterministic 100-pair episode. All policies satisfy farmland-count conservation. Region A shows bimodal behaviour (seed 4 in both configurations fails with $+5\%$ slope while other seeds cluster near zero). Region B's 10 rows all show positive slope change, illustrating the uniform paired-inference failure discussed in Section~\ref{sec:results-mismatch}. Region C is the only region in which most DRL seeds achieve negative slope change.}
\label{tab:per-seed}
\begin{tabular}{llrrrrrrr}
\toprule
Region & Cfg & Seed & Init~$\bar{s}$ & Final~$\bar{s}$ & $\Delta$\,(\%) & Init~$\bar{c}$ & Final~$\bar{c}$ & $\Delta\bar{c}$ \\
\midrule
A  & Lagr & 0 & 8.598 & 8.634 & $+0.41$ & 3.736 & 3.839 & $+0.103$ \\
   & Lagr & 1 & 8.598 & 8.565 & $-0.39$ & 3.736 & 3.880 & $+0.143$ \\
   & Lagr & 2 & 8.598 & 8.596 & $-0.03$ & 3.736 & 3.885 & $+0.149$ \\
   & Lagr & 3 & 8.598 & 8.587 & $-0.13$ & 3.736 & 3.897 & $+0.161$ \\
   & Lagr & 4 & 8.598 & 9.050 & $+5.25$ & 3.736 & 3.738 & $+0.002$ \\
\cmidrule{2-9}
   & Ent  & 0 & 8.598 & 8.608 & $+0.12$ & 3.736 & 3.853 & $+0.117$ \\
   & Ent  & 1 & 8.598 & 8.548 & $-0.59$ & 3.736 & 3.809 & $+0.073$ \\
   & Ent  & 2 & 8.598 & 8.551 & $-0.55$ & 3.736 & 3.864 & $+0.127$ \\
   & Ent  & 3 & 8.598 & 8.515 & $-0.97$ & 3.736 & 3.775 & $+0.039$ \\
   & Ent  & 4 & 8.598 & 9.042 & $+5.17$ & 3.736 & 3.713 & $-0.023$ \\
\midrule
B  & Lagr & 0 & 8.034 & 8.219 & $+2.30$ & 2.720 & 2.553 & $-0.167$ \\
   & Lagr & 1 & 8.034 & 8.187 & $+1.90$ & 2.720 & 2.559 & $-0.161$ \\
   & Lagr & 2 & 8.034 & 8.204 & $+2.12$ & 2.720 & 2.560 & $-0.160$ \\
   & Lagr & 3 & 8.034 & 8.222 & $+2.34$ & 2.720 & 2.554 & $-0.166$ \\
   & Lagr & 4 & 8.034 & 8.218 & $+2.29$ & 2.720 & 2.554 & $-0.166$ \\
\cmidrule{2-9}
   & Ent  & 0 & 8.034 & 8.223 & $+2.36$ & 2.720 & 2.558 & $-0.162$ \\
   & Ent  & 1 & 8.034 & 8.161 & $+1.57$ & 2.720 & 2.563 & $-0.157$ \\
   & Ent  & 2 & 8.034 & 8.226 & $+2.39$ & 2.720 & 2.557 & $-0.163$ \\
   & Ent  & 3 & 8.034 & 8.213 & $+2.23$ & 2.720 & 2.551 & $-0.169$ \\
   & Ent  & 4 & 8.034 & 8.214 & $+2.24$ & 2.720 & 2.739 & $+0.019$ \\
\midrule
C  & Lagr & 0 & 7.234 & 7.181 & $-0.72$ & 4.440 & 4.509 & $+0.068$ \\
   & Lagr & 1 & 7.234 & 7.017 & $-3.00$ & 4.440 & 4.503 & $+0.063$ \\
   & Lagr & 2 & 7.234 & 7.075 & $-2.20$ & 4.440 & 4.442 & $+0.001$ \\
   & Lagr & 3 & 7.234 & 7.052 & $-2.51$ & 4.440 & 4.519 & $+0.078$ \\
   & Lagr & 4 & 7.234 & 7.047 & $-2.58$ & 4.440 & 4.469 & $+0.029$ \\
\cmidrule{2-9}
   & Ent  & 0 & 7.234 & 7.138 & $-1.33$ & 4.440 & 4.502 & $+0.062$ \\
   & Ent  & 1 & 7.234 & 6.993 & $-3.33$ & 4.440 & 4.477 & $+0.037$ \\
   & Ent  & 2 & 7.234 & 7.111 & $-1.70$ & 4.440 & 4.512 & $+0.072$ \\
   & Ent  & 3 & 7.234 & 7.297 & $+0.88$ & 4.440 & 4.560 & $+0.120$ \\
   & Ent  & 4 & 7.234 & 7.094 & $-1.93$ & 4.440 & 4.490 & $+0.050$ \\
\bottomrule
\end{tabular}
\end{table}

\section{Computational reproducibility notes}\label{app:repro}

The 30 DRL training runs were executed on Google Colab Pro+ (NVIDIA A100, 24-hour session limit). Total compute was approximately 75 A100-hours, distributed across multiple sessions. The sweep driver \texttt{colab\_train\_v2.py} writes \texttt{sweep\_state.json} after every completed run, supporting transparent restart across sessions.

\begin{table}[htbp]
\centering
\caption{Training-run audit for the final DRL evaluation set. Failed-write runs were detected before analysis, retrained, and excluded from the final evaluation until valid model artefacts were present.}
\label{tab:compute-audit}
\small
\begin{tabular}{@{}p{0.72\textwidth}r@{}}
\toprule
Audit item & Count/status \\
\midrule
Planned DRL runs in final factorial design & 30 \\
Silent failed Drive writes detected in initial sweep & 8 \\
Affected runs re-queued and retrained & 8 \\
Region B policies retrained under area-weighted environment & 10 \\
Final trained policies included in reported tables & 30 \\
Final deterministic episodes completing the 100-pair budget & 30/30 \\
\bottomrule
\end{tabular}
\end{table}

During execution, a Google Drive quota constraint caused 8 of the initial 30 runs to silently fail-write their model artefacts to Drive while subprocess return codes reported success. We detected this via a post-hoc audit (\texttt{audit\_sweep.py}) that verifies file existence and minimum-byte-size thresholds for every claimed-completed run; the 8 affected runs were re-queued and retrained. The release archive includes both the original and corrected \texttt{sweep\_state.json} for full transparency.

A robust version of the sweep driver, included as \texttt{cell14\_robust.py} in the release, performs local-disk staging followed by byte-verified copy-to-Drive before marking any run complete; this should be used in future replications to avoid silent-failure modes under quota constraints.

The CEUS figure set was generated from the same area-weighted tables used in the manuscript. The figure-generation script validates all 30 trajectory JSON files against the per-seed CSV before plotting, then records the source files in a figure audit JSON released with the review artefacts. This check is intended to make the figure--table correspondence reproducible.

\clearpage

\bibliographystyle{tfcad}
\begin{thebibliography}{99}

\bibitem[Aerts et~al.(2003)]{aerts2003using}
Aerts, J.~C., Eisinger, E., Heuvelink, G.~B., and Stewart, T.~J. (2003).
Using linear integer programming for multi-site land-use allocation.
\textit{Geographical Analysis} 35(2): 148--169.

\bibitem[Achanta et~al.(2012)]{achanta2012slic}
Achanta, R., Shaji, A., Smith, K., Lucchi, A., Fua, P., and S{\"u}sstrunk, S. (2012).
SLIC superpixels compared to state-of-the-art superpixel methods.
\textit{IEEE Transactions on Pattern Analysis and Machine Intelligence} 34(11): 2274--2282.

\bibitem[Bello et~al.(2017)]{bello2017neural}
Bello, I., Pham, H., Le, Q.~V., Norouzi, M., and Bengio, S. (2017).
Neural combinatorial optimization with reinforcement learning.
\textit{International Conference on Learning Representations Workshop Track}.
arXiv:1611.09940.

\bibitem[Cao et~al.(2011)]{cao2011sustainable}
Cao, K., Batty, M., Huang, B., Liu, Y., Yu, L., and Chen, J. (2011).
Spatial multi-objective land use optimization: extensions to the non-dominated sorting genetic algorithm-II.
\textit{International Journal of Geographical Information Science} 25(12): 1949--1969.

\bibitem[Cao et~al.(2012)]{cao2012spatial}
Cao, K., Huang, B., Wang, S., and Lin, H. (2012).
Sustainable land use optimization using boundary-based fast genetic algorithm.
\textit{Computers, Environment and Urban Systems} 36(3): 257--269.

\bibitem[Deb et~al.(2002)]{deb2002fast}
Deb, K., Pratap, A., Agarwal, S., and Meyarivan, T. (2002).
A fast and elitist multiobjective genetic algorithm: NSGA-II.
\textit{IEEE Transactions on Evolutionary Computation} 6(2): 182--197.

\bibitem[Huang and Onta\~{n}\'{o}n(2022)]{huang2022closer}
Huang, S. and Onta\~{n}\'{o}n, S. (2022).
A closer look at invalid action masking in policy gradient algorithms.
\textit{The International FLAIRS Conference Proceedings} 35.
\url{https://doi.org/10.32473/flairs.v35i.130584}.

\bibitem[Karra et~al.(2021)]{karra2021global}
Karra, K., Kontgis, C., Statman-Weil, Z., Mazzariello, J.~C., Mathis, M., and Brumby, S.~P. (2021).
Global land use/land cover with {Sentinel}-2 and deep learning.
In \textit{2021 IEEE International Geoscience and Remote Sensing Symposium}, 4704--4707.
\url{https://doi.org/10.1109/IGARSS47720.2021.9553499}.

\bibitem[Kool et~al.(2019)]{kool2019attention}
Kool, W., van Hoof, H., and Welling, M. (2019).
Attention, learn to solve routing problems!
\textit{International Conference on Learning Representations}.

\bibitem[Liu et~al.(2014)]{liu2014land}
Liu, Y., Fang, F., and Li, Y. (2014).
Key issues of land use in China and implications for policy making.
\textit{Land Use Policy} 40: 6--12.

\bibitem[Long et~al.(2020)]{long2018land}
Long, H., Qu, Y., Tu, S., Zhang, Y., and Jiang, Y. (2020).
Development of land use transitions research in {China}.
\textit{Journal of Geographical Sciences} 30(7): 1195--1214.
\url{https://doi.org/10.1007/s11442-020-1777-9}.

\bibitem[Nazari et~al.(2018)]{nazari2018reinforcement}
Nazari, M., Oroojlooy, A., Snyder, L.~V., and Tak\'{a}\v{c}, M. (2018).
Reinforcement learning for solving the vehicle routing problem.
\textit{Advances in Neural Information Processing Systems}, 9861--9871.

\bibitem[Niu et~al.(2023)]{niu2023ffi}
Niu, L., Xu, M., and de Vries, W.~T. (2023).
Investigation and comparison of spatial--temporal characteristics of farmland fragmentation in the {Beijing}--{Tianjin}--{Hebei} region, {China}, and {Bavaria}, {Germany}.
\textit{Land} 12(7): 1410.
\url{https://doi.org/10.3390/land12071410}.

\bibitem[Stewart et~al.(2004)]{stewart2004multiobjective}
Stewart, T.~J., Janssen, R., and van Herwijnen, M. (2004).
A genetic algorithm approach to multiobjective land use planning.
\textit{Computers \& Operations Research} 31(14): 2293--2313.

\bibitem[Stooke et~al.(2020)]{stooke2020responsive}
Stooke, A., Achiam, J., and Abbeel, P. (2020).
Responsive safety in reinforcement learning by PID Lagrangian methods.
\textit{International Conference on Machine Learning}, 9133--9143.

\bibitem[Tan et~al.(2006)]{tan2006fragmentation}
Tan, S., Heerink, N., and Qu, F. (2006).
Land fragmentation and its driving forces in China.
\textit{Land Use Policy} 23(3): 272--285.

\bibitem[Tessler et~al.(2019)]{tessler2019reward}
Tessler, C., Mankowitz, D.~J., and Mannor, S. (2019).
Reward constrained policy optimization.
\textit{International Conference on Learning Representations}.

\bibitem[Wu et~al.(2005)]{wu2005consolidation}
Wu, Z., Liu, M., and Davis, J. (2005).
Land consolidation and productivity in Chinese household crop production.
\textit{China Economic Review} 16(1): 28--49.

\bibitem[Zheng et~al.(2023)]{zheng2023spatial}
Zheng, Y., Lin, Y., Zhao, L., Wu, T., Jin, D., and Li, Y. (2023).
Spatial planning of urban communities via deep reinforcement learning.
\textit{Nature Computational Science} 3: 748--762.

\bibitem[Schulman et~al.(2017)]{schulman2017proximal}
Schulman, J., Wolski, F., Dhariwal, P., Radford, A., and Klimov, O. (2017).
Proximal policy optimization algorithms.
arXiv:1707.06347.

\bibitem[Raffin et~al.(2021)]{raffin2021stablebaselines}
Raffin, A., Hill, A., Gleave, A., Kanervisto, A., Ernestus, M., and Dormann, N. (2021).
Stable-Baselines3: reliable reinforcement learning implementations.
\textit{Journal of Machine Learning Research} 22(268): 1--8.

\bibitem[Rey and Anselin(2007)]{rey2007pysal}
Rey, S.~J. and Anselin, L. (2007).
PySAL: a Python library of spatial analytical methods.
\textit{The Review of Regional Studies} 37(1): 5--27.
\url{https://doi.org/10.52324/001c.8285}.

\bibitem[European Space Agency(2021)]{esa2021copernicusdem}
European Space Agency (2021).
Copernicus Global Digital Elevation Model GLO-30.
\url{https://spacedata.copernicus.eu/web/cscda/dataset-details?articleId=394198}.

\bibitem[Horn(1981)]{horn1981slope}
Horn, B.~K.~P. (1981).
Hill shading and the reflectance map.
\textit{Proceedings of the IEEE} 69(1): 14--47.
\url{https://doi.org/10.1109/PROC.1981.11918}.

\end{thebibliography}

\end{document}
